% Fate's Edge SRD −−− Summoning (Section 13)

\section{Summoning}

Summoning is the art of calling beings from beyond the mortal veil — spirits, elementals, the dead, fey, demons, or stranger entities. It is always dangerous, often costly, and never without consequence.

\subsection{Who Can Summon?}

Any character may attempt a summoning ritual, but only those with the \textbf{Pact−Whisperer} talent can reliably bind and command spirits in combat time. NPCs and plot events use the same principles: a summoner prepares a circle, pays a cost, and risks backlash.

\subsection{Summoning for Narrative (Greater Powers)}

Not all summonings are tactical. Sometimes the story demands calling a greater power — a patron−level entity, a forgotten god, or a primordial force. These rituals are not resolved with a single roll; they are extended challenges that drive the plot.

⫕section{Ritual Structure for Greater Summonings}

\begin{itemize}
  ℹtem \textbf{Preparation Timer [6–12]:} The summoner gathers rare components, secures a threshold, and learns the entity's true name (or equivalent).
  ℹtem \textbf{The Calling:} A dramatic scene where the summoner rolls \texttt{Spirit + Lore} vs. DV (set by the entity's power). Failure may summon something else.
  ℹtem \textbf{The Bargain:} Greater powers do not serve for free. They demand a price — a memory, a firstborn, a deed, or an ongoing Obligation (tracked as a campaign front).
  ℹtem \textbf{Consequences:} On success, the entity manifests (briefly) and may grant a boon, answer a question, or perform a single impossible task. On failure, the entity corrupts, attacks, or curses the summoner.
\end{itemize}

\textbf{Example: Summoning the Pale Shepherd's Herald}
A village elder attempts to summon a death−herald to save a dying child. Preparation timer [8]: gather grave−ash, a silver bell, and the child's true name. Calling roll: Spirit 2 + Lore 3 = 5 dice vs DV 5 $→$ Partial. The herald appears but demands a year of the elder's remaining life. The elder accepts; the child lives.

\subsection{PC Summoner (Pact−Whisperer)}

Player characters who specialize in summoning take the \textbf{Pact−Whisperer} talent (2 XP). This allows them to call and bind spirits for temporary aid in scenes.

⫕section{Summoning Talents (Summary)}

Refer to Section~\ref{sec:talents} for full talent descriptions. The essential summoning talents are:
\begin{itemize}
  ℹtem \textbf{Pact−Whisperer (2 XP):} Base access to summoning.
  ℹtem \textbf{Lesser Pactwright (2 XP):} Call spirits of Cap 1.
  ℹtem \textbf{Greater Pactwright (4 XP):} Call spirits of Cap 3.
  ℹtem \textbf{Spirit Link (10 XP):} Spirits act on your turn; all commands are free actions; reduce Leash ticking for natural behaviours by 1.
\end{itemize}

\subsection{Core Mechanics – Call, Bind, and the Leash}

Summoning a spirit follows three steps: Call, Bind, and Leash management.

⫕section{1. Call (1 Action)}
The spirit manifests at Near range. Choose a spirit template (see sample templates below).

⫕section{2. Bind (1 Boon or 1 Fatigue)}
Spend 1 Boon or mark 1 Fatigue to establish control. The spirit is now bound and will follow simple commands.

⫕section{3. Leash Capacity}
The spirit has a \textbf{Leash} — a track of segments equal to \texttt{Cap + Spirit}. Cap is the spirit's tier: 1 for Lesser, 3 for Greater. The Leash represents the strain of keeping the spirit in the mortal world.

⫕section{Ticking the Leash}
The Leash ticks forward (gains segments) under the following conditions:
\begin{itemize}
  ℹtem The spirit takes Harm (1 segment per level of Harm).
  ℹtem You command the spirit to act against its nature (1 segment).
  ℹtem You split focus — take another significant action while the spirit is acting (1 segment).
  ℹtem The spirit crosses a [WARD] without dispelling it (DV = Cap; tick segments equal to the ward's DV).
  ℹtem (Optional narrative trigger) A rival contests control, or the spirit moves too far from you.
\end{itemize}

⫕section{When the Leash Fills}
When the Leash reaches its capacity, the spirit acts once according to its nature (the GM decides the action, which is often destructive or self−serving), then departs immediately (or turns hostile at the GM's discretion). The bond is broken.

⫕section{Boon Finesse}
Once per round, you may spend \textbf{1 Boon} to clear 1 tick from the current spirit's Leash (cannot be used after the Leash has already filled). This represents appeasing the spirit, reaffirming control, or offering a small token.

\subsection{Commanding Spirits}

\begin{itemize}
  ℹtem \textbf{Quick Command:} Simple orders (``attack that enemy'', ``move there'', ``defend me'', ``retrieve that object'') do not consume your Action. You may issue one Quick Command per round as a free action.
  ℹtem \textbf{Complex Command:} If the order requires interpretation, multiple steps, or finesse (``guard this door and alert me if anyone passes'', ``search the room for a hidden key''), it uses your Action for the turn.
  ℹtem \textbf{Spirit Link (Major Talent):} Spirits act on your turn; all commands are free actions; reduce Leash ticking for natural behaviours by 1.
\end{itemize}

\subsection{Spirit Bond Progression}

For spirits you summon regularly (e.g., a favoured elemental or a recurring shade), track a \textbf{Spirit Bond Timer [4]}.

\begin{itemize}
  ℹtem Mark segments for successful commands, shared victories, or mutual aid (once per type per session).
  ℹtem At \textbf{2 segments:} Gain +1 die to communicate with this spirit type.
  ℹtem At \textbf{4 segments:} The spirit becomes \textbf{Favored}. Choose one: Leash Capacity +1, or –1 tick per command (minimum 0). Additionally, the spirit grants you +1 Boon when it departs naturally (not destroyed or forced out).
\end{itemize}

\subsection{Sample Spirit Templates}

⫕section{Lesser (Cap 1)}

\begin{itemize}
  ℹtem \textbf{Shade:} Stealthy scout; can pass through cracks; Harm 1, Leash = 1+Spirit.
  ℹtem \textbf{Emberling:} Small fire elemental; deals 1 Harm (fire) on touch; can ignite objects.
  ℹtem \textbf{Tangle−root:} Plant spirit; can entangle one target (Entangled condition, Athletics DV 3 to break).
\end{itemize}

⫕section{Greater (Cap 3)}

\begin{itemize}
  ℹtem \textbf{Flame Elemental:} Area fire attack (Harm 2 to all in Close); immune to mundane fire; Leash = 3+Spirit.
  ℹtem \textbf{Shadow Hound:} Track any scent; can move silently; attacks with Harm 2 (shadow bite).
  ℹtem \textbf{Stone Guardian:} High defense; can interpose and reduce Harm to an ally by 1 (once per round). Harm 1, but +2 Armor.
\end{itemize}

\subsection{Specializations (for recurring spirits)}

Choose one specialization when a spirit becomes \textbf{Favored}:
\begin{itemize}
  ℹtem \textbf{Combat Specialist:} +1 Harm on attacks; ignores first Harm 1 per scene.
  ℹtem \textbf{Scout Form:} Stealthy; Cap 1 can carry 2 kg, Cap 3 can carry 10 kg; +2 dice to Stealth.
  ℹtem \textbf{Utility Spirit:} Can perform simple tasks (opening unlocked doors, carrying light objects, fetching items).
  ℹtem \textbf{Shield Guardian:} Once per round, interpose and convert an ally's Harm 1 into Fatigue on the spirit.
\end{itemize}

\subsection{GM Guidelines for Summoning}

⫕section{NPC Summoners}
Treat enemy summoners as having the same limitations: they must Call (action) and Bind (Boon or Fatigue). Their spirits have a Leash timer. The GM may choose not to track Leash for minor NPCs if the fiction does not demand it, but for major antagonists, use the same mechanics — this creates tactical choices (e.g., attack the summoner to break control, or force the spirit to cross a ward).

⫕section{Plot−Level Summonings}
When the story requires a massive summoning (a demon army, a dead god, the Hollow Bride), do not use PC mechanics. Instead:
\begin{itemize}
  ℹtem Create a \textbf{Ritual Timer [6–12]} that the party can disrupt.
  ℹtem Each segment represents a component (blood sacrifice, alignment of stars, chanting cultists).
  ℹtem The party can sabotage components, kill cultists, or break the circle to halt the summoning.
  ℹtem If the timer fills, the entity arrives — now the party must bargain, fight, or flee.
\end{itemize}

⫕section{Hazards of Greater Summoning}
\begin{itemize}
  ℹtem \textbf{Uncontrolled Manifestation:} The entity may be hostile to everyone, including the summoner.
  ℹtem \textbf{Collateral Damage:} Reality warps, wards fail, the dead rise, or the environment becomes a hazard.
  ℹtem \textbf{Debt:} Even if bound, the entity extracts a price — an Obligation timer that ticks per scene, demanding sacrifice or service.
\end{itemize}

\textbf{Example: The Cult of the Ninth Bell}
Cultists attempt to summon a fragment of the Ninth (void entity). A 6−segment timer: collect nine souls (segment 1), ring the forbidden bell (segment 2), etc. The party can interrupt at any segment. If the timer fills, the Ninth Rim appears — treat as a Tier IV hazard with a Bargain (offer a memory) or a Catastrophe (reality starts erasing the town).

